%% ma: L11-B20-0002
%% lop: 11
%% chuong: 6
%% bai: 20
%% dang: 11.20.2
%% loai: TN4
%% muc: TH
%% dap_an: "D"
%% nguon: "(NBV)-ĐỀ SỐ 15-MỖI NGÀY 1 ĐỀ THI 2025.pdf (D15, MỖI NGÀY 1 ĐỀ - Đề số 15), Phần I Câu 2"
%% kien_thuc: "Bài 20 (hàm số lôgarit: tính chất, tập giá trị, đồ thị)"
%% trang_thai: chinh-thuc
%% ghi_chu: "Đã duyệt 16/10/2026. Đáp án gốc D đúng."
\begin{ex}
Cho hàm số $y=\log_{0{,}5}x$. Mệnh đề nào sau đây sai?
\choice{Hàm số đã cho nghịch biến trên $(0;+\infty)$}{Đồ thị hàm số đã cho nằm phía bên phải trục tung}{Đồ thị hàm số đã cho đi qua điểm $A(1;0)$}{\True Tập giá trị của hàm số đã cho là $T=(0;+\infty)$}
\loigiai{Cơ số $0{,}5<1$ nên hàm số nghịch biến trên $(0;+\infty)$; đồ thị nằm bên phải trục tung và đi qua $A(1;0)$ vì $\log_{0{,}5}1=0$. Tập giá trị của hàm số là $\mathbb R$, nên mệnh đề D sai.}
\end{ex}
